% Relativistic Chapter XIII, §119: Jeans Gravitational Instability Criterion
% Relativistic generalization of Chandrasekhar, Chapter XIII §119
% Conventions: see RELATIVISTIC_CONVENTIONS.md

\section{Relativistic Jeans gravitational instability criterion}
\label{sec:rel-13-119}

%% ================================================================
\subsection{Classical Jeans instability: a brief recapitulation}
\label{sec:rel-13-119a}

In Chapter~XIII, \S119 of the non-relativistic treatment,
Chandrasekhar derives the gravitational instability of an infinite
homogeneous medium at rest---Jeans's instability.  In the Newtonian
theory the perturbation equations for the velocity $\mathbf{u}$,
the density fluctuation $\delta\rho$, and the gravitational
potential fluctuation $\delta\mathfrak{B}$ are
\begin{equation}
\label{eq:rel-13-79}
  \rho\frac{\partial\mathbf{u}}{\partial t}
  = -\operatorname{grad}\delta p
    + \rho\,\operatorname{grad}\delta\mathfrak{B},
\end{equation}
\begin{equation}
\label{eq:rel-13-80}
  \frac{\partial}{\partial t}\delta\rho
  = -\rho\,\operatorname{div}\mathbf{u},
\end{equation}
\begin{equation}
\label{eq:rel-13-81}
  \nabla^{2}\delta\mathfrak{B} = -4\pi G\,\delta\rho.
\end{equation}
With the adiabatic relation $\delta p = \cs^{2}\,\delta\rho$
(where $\cs^{2}=\gamma p/\rho$ is the squared sound speed), one
obtains the dispersion relation
\begin{equation}
\label{eq:rel-13-87-cl}
  \omega^{2} = \cs^{2}\,k^{2} - 4\pi G\rho.
\end{equation}
The Jeans wavenumber $k_{J}$ and Jeans length $\lambda_{J}$ are
\begin{equation}
\label{eq:jeans-classical}
  k_{J}^{2} = \frac{4\pi G\rho}{\cs^{2}},
  \qquad
  \lambda_{J}
  = \frac{2\pi}{k_{J}}
  = \cs\sqrt{\frac{\pi}{G\rho}}\,.
\end{equation}
Modes with $k < k_{J}$ (wavelengths exceeding $\lambda_{J}$)
are gravitationally unstable: the self-gravity of the density
perturbation overwhelms the restoring pressure force and the
perturbation grows exponentially.

%% ================================================================
\subsection{Why a relativistic treatment is necessary}
\label{sec:rel-13-119b}

In general relativity, the source of gravity is not merely mass
density but energy density---and, crucially, \emph{pressure itself
gravitates}.  This has two profound consequences for Jeans
instability:

\begin{enumerate}
\item \textbf{Inertia of energy.}  The gravitating mass density
  is replaced by the relativistic enthalpy density
  $\enthalpy = (\varepsilon + p)/c^{2}$, where $\varepsilon$ is
  the total energy density (including rest-mass energy) and $p$ is
  the pressure.  This enhances the gravitational source for any
  medium in which the pressure is non-negligible compared with the
  energy density (e.g.\ radiation-dominated plasmas, neutron-star
  interiors, the early universe).

\item \textbf{Pressure as a source of gravity.}  In the Newtonian
  theory, pressure enters only as a \emph{restoring} force opposing
  gravitational collapse.  In general relativity, pressure also
  appears in the gravitational source term.  The Tolman--Oppenheimer
  --Volkoff equation, for instance, contains the combination
  $(\varepsilon + 3p)/(2c^{2})$ as the active gravitational mass
  density, reflecting the trace of the energy-momentum tensor.
  When the pressure is large, gravity is \emph{enhanced} and the
  medium becomes more susceptible to collapse---the opposite of the
  na\"{\i}ve expectation from Newtonian intuition.
\end{enumerate}

This is the key Chandrasekhar insight: \textbf{in general
relativity, pressure destabilises rather than stabilises}.  At
sufficiently high densities, no amount of pressure can prevent
gravitational collapse---the physics underlying the Chandrasekhar
mass limit for white dwarfs and the Tolman--Oppenheimer--Volkoff
limit for neutron stars.

Relativistic Jeans instability is directly relevant to:
\begin{itemize}
\item structure formation in the radiation-dominated early
  universe, where $p = \varepsilon/3$;
\item neutron-star interiors, where $p/\varepsilon \sim 0.1$--$0.3$;
\item quark--gluon plasmas produced in heavy-ion collisions.
\end{itemize}

%% ================================================================
\subsection{Relativistic perturbation equations}
\label{sec:rel-13-119c}

We work in a flat Friedmann-like background (Minkowski spacetime
perturbed by weak self-gravity) with signature $(-,+,+,+)$ and
keep $c$ explicit throughout.  The background is a homogeneous
perfect fluid at rest with energy density $\varepsilon$, pressure
$p$, and four-velocity $u^{\mu} = (c, 0, 0, 0)$.  The
energy-momentum tensor is
\begin{equation}
\label{eq:rel-13-Tmunu}
  T^{\mu\nu}
  = \frac{\varepsilon + p}{c^{2}}\,u^{\mu}u^{\nu}
    + p\,g^{\mu\nu}.
\end{equation}

In the weak-field, slow-motion limit the relativistic analogue of
Poisson's equation~\eqref{eq:rel-13-81} is obtained from the
linearised Einstein equations.  The $00$-component of the
Einstein equation gives
\begin{equation}
\label{eq:rel-13-poisson-rel}
  \nabla^{2}\delta\Phi
  = -4\pi G\,\frac{\delta\varepsilon + 3\,\delta p}{2\,c^{2}}\,,
\end{equation}
where $\delta\Phi$ is the Newtonian-limit gravitational potential.
For a perfect fluid with adiabatic perturbations $\delta p =
\cs^{2}\,\delta\varepsilon/c^{2}$ (since $\cs^{2} = (\partial
p/\partial\varepsilon)_{s}\,c^{2}$ in the relativistic convention
of~RELATIVISTIC\_CONVENTIONS.md), this becomes
\begin{equation}
\label{eq:rel-13-poisson-rel2}
  \nabla^{2}\delta\Phi
  = -\frac{4\pi G}{2\,c^{2}}
    \bigl(1 + 3\,\cs^{2}/c^{2}\bigr)\,\delta\varepsilon\,.
\end{equation}

Simultaneously, the relativistic Euler equation for the
perturbations reads (in the rest frame of the background fluid)
\begin{equation}
\label{eq:rel-13-euler-rel}
  \frac{\varepsilon + p}{c^{2}}\,
  \frac{\partial\mathbf{u}}{\partial t}
  = -\operatorname{grad}\delta p
    + \frac{\varepsilon + p}{c^{2}}\,
      \operatorname{grad}\delta\Phi\,,
\end{equation}
and the energy conservation equation is
\begin{equation}
\label{eq:rel-13-cont-rel}
  \frac{\partial}{\partial t}\delta\varepsilon
  = -(\varepsilon + p)\,\operatorname{div}\mathbf{u}\,.
\end{equation}
Note that the inertia in eq.~\eqref{eq:rel-13-euler-rel} is the
enthalpy density $\enthalpy = (\varepsilon + p)/c^{2}$, not the
rest-mass density $\rho$.  Likewise, eq.~\eqref{eq:rel-13-cont-rel}
involves $(\varepsilon + p)$ rather than $\rho$ in the source term.

%% ================================================================
\subsection{Relativistic dispersion relation}
\label{sec:rel-13-119d}

Taking the divergence of eq.~\eqref{eq:rel-13-euler-rel},
substituting eq.~\eqref{eq:rel-13-cont-rel} and
eq.~\eqref{eq:rel-13-poisson-rel2}, and seeking plane-wave
solutions $\delta\varepsilon \propto
e^{i(\mathbf{k}\cdot\mathbf{x} - \omega t)}$, we obtain the
\textbf{relativistic Jeans dispersion relation}:
\begin{equation}
\label{eq:rel-13-dispersion}
\boxed{
  \omega^{2}
  = \cs^{2}\,k^{2}
    - \frac{4\pi G}{c^{2}}\,
      (\varepsilon + p)\,
      \frac{1 + 3\,\cs^{2}/c^{2}}{2}\,.
}
\end{equation}

To make the structure transparent, define the relativistic
enthalpy density $\enthalpy \equiv (\varepsilon + p)/c^{2}$ and
the ``active gravitational factor''
$\mathcal{A} \equiv (1 + 3\,\cs^{2}/c^{2})/2$.  Then:
\begin{equation}
\label{eq:rel-13-dispersion-compact}
  \omega^{2}
  = \cs^{2}\,k^{2}
    - 4\pi G\,\enthalpy\,\mathcal{A}\,.
\end{equation}
In the Newtonian limit $p \ll \varepsilon \approx \rho c^{2}$,
$\cs^{2} \ll c^{2}$, we recover the classical result~\eqref{eq:rel-13-87-cl}:
\begin{equation}
  \omega^{2}
  \;\xrightarrow{\;c\to\infty\;}
  \;\cs^{2}\,k^{2} - 4\pi G\rho\,.
\end{equation}

%% ================================================================
\subsection{Relativistic Jeans wavenumber, length, and mass}
\label{sec:rel-13-119e}

Instability ($\omega^{2} < 0$) occurs for all wavenumbers
$k < k_{J,\mathrm{rel}}$, where
\begin{equation}
\label{eq:rel-13-kJ}
\boxed{
  k_{J,\mathrm{rel}}^{2}
  = \frac{4\pi G}{c^{2}\,\cs^{2}}\,
    (\varepsilon + p)\,
    \frac{1 + 3\,\cs^{2}/c^{2}}{2}\,.
}
\end{equation}
Equivalently,
\begin{equation}
\label{eq:rel-13-kJ-expanded}
  k_{J,\mathrm{rel}}^{2}
  = \frac{4\pi G\,\enthalpy}{\cs^{2}}\,
    \Bigl(
      1 + \frac{3\,\cs^{2}}{2\,c^{2}}
      + \mathcal{O}\bigl(\cs^{4}/c^{4}\bigr)
    \Bigr).
\end{equation}
Comparing with the classical expression
$k_{J}^{2} = 4\pi G\rho/\cs^{2}$, we see two relativistic
enhancements:
\begin{enumerate}
\item $\rho \to \enthalpy = (\varepsilon + p)/c^{2}$: the inertia
  of internal energy and pressure increases the effective
  gravitating mass;
\item the factor $\mathcal{A} = (1 + 3\cs^{2}/c^{2})/2 > 1/2$:
  pressure enters the gravitational source term,
  \emph{further enhancing the effective gravity}.
\end{enumerate}
Both effects act in the same direction: they
\textbf{increase} $k_{J}$, meaning instability sets in at
\textbf{shorter} wavelengths than Newtonian theory predicts.

The relativistic Jeans length and Jeans mass are
\begin{equation}
\label{eq:rel-13-lambdaJ}
  \lambda_{J,\mathrm{rel}}
  = \frac{2\pi}{k_{J,\mathrm{rel}}}
  = \cs\,
    \sqrt{
      \frac{\pi\,c^{2}}
           {G\,(\varepsilon + p)\,\mathcal{A}}
    }
  \;<\; \lambda_{J,\mathrm{classical}}\,,
\end{equation}
\begin{equation}
\label{eq:rel-13-MJ}
  M_{J,\mathrm{rel}}
  = \frac{4\pi}{3}\,\enthalpy\,
    \Bigl(\frac{\lambda_{J,\mathrm{rel}}}{2}\Bigr)^{\!3}
  \;<\; M_{J,\mathrm{classical}}\,.
\end{equation}
The inequalities hold because the relativistic corrections to
gravity are always positive for physically reasonable equations of
state ($\varepsilon > 0$, $p \geq 0$).

\paragraph{Key physical consequence.}
In general relativity the Jeans mass is \emph{smaller} than its
Newtonian counterpart.  Put differently, lumps of matter that
would be Jeans-stable in Newtonian gravity are already unstable
when relativistic gravity is properly accounted for.  This is the
\emph{pressure-destabilisation} effect: the very pressure that
supports the medium against collapse simultaneously contributes to
the gravitational source, weakening the net restoring force.

%% ================================================================
\subsection{The pressure-destabilisation mechanism in detail}
\label{sec:rel-13-119f}

To see why pressure aids collapse in general relativity, consider
the energy argument of Chandrasekhar.  A perturbation of
wavenumber $k$ involves:
\begin{itemize}
\item \textbf{Pressure work (restoring):} proportional to
  $\cs^{2}\,k^{2}$, independent of $G$.
\item \textbf{Gravitational energy release (destabilising):}
  proportional to $4\pi G\,\enthalpy\,\mathcal{A}/c^{2}$.
\end{itemize}
In the Newtonian limit $\mathcal{A} \to 1/2$ and $\enthalpy \to \rho$,
so the gravitational term is $\sim 4\pi G\rho$, and only the
\emph{mass} contributes to gravity.  In the relativistic case:
\begin{equation}
  4\pi G\,\enthalpy\,\mathcal{A}
  = \frac{4\pi G}{c^{2}}\,
    \frac{(\varepsilon + p)(1 + 3\cs^{2}/c^{2})}{2}\,,
\end{equation}
where the pressure appears \emph{both} in $\enthalpy$ (inertial
contribution) \emph{and} in $\mathcal{A}$ (active gravitational
contribution).  When $p$ is increased at fixed $\varepsilon$:
\begin{enumerate}
\item the restoring force increases as $\cs^{2}k^{2}$, but only
  linearly in the sound speed;
\item the gravitational source increases through both $\enthalpy$
  and $\mathcal{A}$, and for a stiff equation of state
  ($\cs^{2} \lesssim c^{2}$) this increase can dominate.
\end{enumerate}
The net result is that the critical wavenumber $k_{J,\mathrm{rel}}$
grows with increasing pressure-to-energy ratio $p/\varepsilon$,
and the Jeans mass decreases.  This is the mechanism behind the
Chandrasekhar and Tolman--Oppenheimer--Volkoff mass limits.

%% ================================================================
\subsection{Causality constraints}
\label{sec:rel-13-119g}

The relativistic treatment imposes strict constraints on
perturbation propagation that have no classical analogue:

\begin{enumerate}
\item \textbf{Perturbation speed.}  Sound waves---the restoring
  mechanism against gravitational collapse---propagate at
  $\cs$, which must satisfy the \emph{causality bound}
  \begin{equation}
  \label{eq:rel-13-causality-cs}
    \cs^{2}
    = \biggl(\frac{\partial p}{\partial\varepsilon}\biggr)_{\!s}
    \leq c^{2}\,.
  \end{equation}
  An equation of state violating this bound would permit
  superluminal signal propagation and is unphysical.  This places
  an upper limit on the stiffness of matter and hence on the
  maximum sound speed available for gravitational support.

\item \textbf{Gravitational response.}  In general relativity,
  changes in the gravitational field propagate at the speed of
  light $c$, not instantaneously as in Newtonian gravity.  The
  gravitational potential perturbation $\delta\Phi$ satisfies a
  wave equation with speed $c$.  In the quasi-static limit used
  in the derivation above (wavelengths much smaller than the
  Hubble radius), retardation effects are negligible and the
  Poisson-equation form~\eqref{eq:rel-13-poisson-rel} is valid.
  For cosmological-scale perturbations, the full linearised
  Einstein equations must be used, introducing additional
  relativistic corrections.

\item \textbf{Maximum sound speed and minimum Jeans mass.}
  Setting $\cs = c$ (the stiffest causal equation of state) in
  eq.~\eqref{eq:rel-13-kJ} gives the \emph{maximum possible}
  Jeans wavenumber:
  \begin{equation}
  \label{eq:rel-13-kJ-max}
    k_{J,\max}^{2}
    = \frac{4\pi G\,(\varepsilon + p)}{c^{4}}\,
      \frac{1 + 3}{2}
    = \frac{8\pi G\,(\varepsilon + p)}{c^{4}}\,,
  \end{equation}
  and hence the minimum Jeans mass below which no causal equation
  of state can stabilise a self-gravitating configuration.
\end{enumerate}

%% ================================================================
\subsection{Limiting cases}
\label{sec:rel-13-119h}

\paragraph{(a) Non-relativistic limit.}
For $p \ll \varepsilon \approx \rho c^{2}$ and
$\cs^{2} \ll c^{2}$:
\begin{equation}
  k_{J,\mathrm{rel}}^{2}
  \to \frac{4\pi G\rho}{\cs^{2}}
  = k_{J}^{2}\,,
  \qquad
  \omega^{2}
  \to \cs^{2}k^{2} - 4\pi G\rho\,,
\end{equation}
recovering the Jeans criterion of \S\ref{sec:rel-13-119a}.

\paragraph{(b) Ultra-relativistic (radiation-dominated) limit.}
For a radiation fluid, $p = \varepsilon/3$ and
$\cs^{2} = c^{2}/3$:
\begin{equation}
\label{eq:rel-13-radiation}
  k_{J,\mathrm{rad}}^{2}
  = \frac{4\pi G}{c^{4}/3}\,
    \frac{4\varepsilon/3 \cdot (1+1)}{2}
  = \frac{4\pi G}{c^{4}/3}\,\frac{4\varepsilon}{3}
  = \frac{16\pi G\varepsilon}{c^{4}}\,,
\end{equation}
and the Jeans length is
$\lambda_{J,\mathrm{rad}} = (c/2)\sqrt{\pi/(G\varepsilon/c^{2})}$.
The numerical factor differs from the Newtonian result by a factor
of order unity, but the important point is that $\varepsilon$
(not $\rho$) is the gravitating source, and the effective gravity
is enhanced by the radiation pressure.

\paragraph{(c) Maximally stiff equation of state.}
For $\cs = c$ (the causal limit, e.g.\ a free scalar field),
with $p = \varepsilon$:
\begin{equation}
  k_{J,\max}^{2}
  = \frac{4\pi G \cdot 2\varepsilon}{c^{2}\cdot c^{2}}
    \cdot \frac{1+3}{2}
  = \frac{16\pi G\varepsilon}{c^{4}}\,.
\end{equation}
This coincides with the radiation result, a reflection of the fact
that both the causal-limit fluid and radiation give the largest
possible gravitational source at given $\varepsilon$.

%% ================================================================
\subsection{Summary of relativistic Jeans criterion}
\label{sec:rel-13-119-summary}

The main results of this section are:

\begin{enumerate}
\item The \textbf{relativistic Jeans dispersion relation}
  is eq.~\eqref{eq:rel-13-dispersion}:
  \[
    \omega^{2}
    = \cs^{2}k^{2}
      - \frac{4\pi G(\varepsilon+p)}{c^{2}}\,
        \frac{1 + 3\cs^{2}/c^{2}}{2}\,.
  \]

\item The \textbf{relativistic Jeans wavenumber}
  is eq.~\eqref{eq:rel-13-kJ}:
  \[
    k_{J,\mathrm{rel}}^{2}
    = \frac{4\pi G(\varepsilon+p)}{\cs^{2}\,c^{2}}\,
      \frac{1 + 3\cs^{2}/c^{2}}{2}\,.
  \]

\item \textbf{Pressure destabilises} in general relativity:
  $M_{J,\mathrm{rel}} < M_{J,\mathrm{classical}}$ because
  pressure contributes to the gravitational source.  This is the
  physical mechanism behind the Chandrasekhar and
  Tolman--Oppenheimer--Volkoff mass limits.

\item \textbf{Causality constraint}: perturbations propagate at
  $\cs \leq c$; the maximum Jeans wavenumber occurs at $\cs = c$
  and sets the minimum mass that can be gravitationally stable.
\end{enumerate}
